\section{Motivation: Connection Entropy}

We use connection entropy to describe the amplification from semantic design intent to primitive RTL connections:

\begin{equation}
  CE = \frac{\text{primitive signal edges}}{\text{semantic interface edges}}.
\end{equation}

A designer may intend a simple chain from a producer through a FIFO to a consumer. In primitive RTL, that chain expands into many individual signal connections. As connection entropy grows, LLMs are more likely to omit ports, reverse ready/valid direction, mismatch payload widths, confuse reset polarity, cross clock domains without synchronization, generate inconsistent parameters, or produce wrappers that simulate only because the testbench misses the actual protocol bug.

Prompt engineering can reduce some failures, but it does not change the underlying problem: free-form RTL text generation leaves too many coupled correctness obligations in one target string. MICO instead constrains the output space. The model proposes a structured design graph, adapter plan, and contract skeleton. The compiler rejects unsafe connections and returns structured diagnostics that can drive a repair prompt.

The design principles are direct:

\begin{itemize}
  \item LLMs produce candidates, not trusted RTL.
  \item Direction, width, protocol, and domain mismatches are compile-time errors.
  \item CDC and RDC crossings require explicit adapters.
  \item Automatically inserted buffers, pipelines, width adapters, or CDC FIFOs must appear in an elaboration report.
  \item Verification collateral is generated from interface contracts.
\end{itemize}

This framing distinguishes two kinds of complexity. \emph{Local RTL complexity} is the logic inside a module: arithmetic, state machines, memories, and datapaths. \emph{Composition complexity} is the set of obligations required to connect already-existing modules safely. MICO focuses on the latter. A direct RTL prompt must solve both at once. A MICO prompt only proposes a semantic graph; the compiler expands that graph into primitive wiring and rejects illegal edges before RTL tools are asked to elaborate it.

\begin{table}[t]
  \centering
  \caption{Typical composition failures targeted by MICO}
  \begin{tabular}{lll}
    \toprule
    Failure & Free-form RTL symptom & MICO gate \\
    \midrule
    Direction & ready/valid reversed & role check \\
    Width & silent truncation & field width check \\
    Protocol & ad hoc handshake & interface contract \\
    CDC & direct cross-domain wire & adapter check \\
    Reset & mixed polarity/domain & clockdom metadata \\
    \bottomrule
  \end{tabular}
\end{table}
